\def\xfrac#1#2{#1/#2}

\def\lfb{\begin{tikzpicture}
    \def\cA{\begin{tikzpicture}\draw[fill=white](0,0) circle(.13);
        \node at(0,0){A};
      \end{tikzpicture}}
    \def\cV{\begin{tikzpicture}\draw[fill=white](0,0) circle(.13);
        \node at(0,-.02){V};
      \end{tikzpicture}}
    \draw(3.5,0)--(0,0)--(0,1.3)--(3.5,1.3);
    \draw(1.3,0)--(1.3,1.3);
    \node at(1.3,.65){\cV};
\draw[fill=black](1.3,0) circle(.03);
\draw[fill=black](1.3,1.3) circle(.03);
    \begin{scope}[xshift=25]
    \draw(1.3,0)--(1.3,1.3);
    \node at(1.3,.65){\cV};
\draw[fill=black](1.3,1.3) circle(.03);
\draw[fill=black](1.3,0) circle(.03);
  \end{scope}
  \node at(.65,1.3){\cA};
  \node at(1.75,1.3){\cA};
  \node at(2.65,1.3){\cA};
\draw[fill=black](3.6,0) circle(.03);
\draw[fill=black](3.7,0) circle(.03);
\draw[fill=black](3.8,0) circle(.03);
\draw[fill=black](3.6,1.3) circle(.03);
\draw[fill=black](3.7,1.3) circle(.03);
\draw[fill=black](3.8,1.3) circle(.03);
\draw[white](0,.63)--(0,.7);
\draw[thick](-.15,.7)--(.15,.7);
\draw[thick](-.1,.63)--(.1,.63);
\end{tikzpicture}}


\def\rysc{\begin{tikzpicture}
  \draw[fill=black](0,0) circle(.03) ;
  \draw[fill=black](0,1.3) circle(.03) ;
  \draw[fill=black](.8,0) circle(.03) ;
  \draw[fill=black](.8,1.3) circle(.03) ;
  \draw[fill=black](3,0) circle(.03) ;
  \draw[fill=black](3,1.3) circle(.03) ;
  \draw(0,0)--(0,1.3);
  \draw(3,0)--(3,1.3);
  \draw(2,0)--(2,1.3);
  \draw[->](.8,1.3)--(1,1.3);
 \draw[->](1,1.3)--(2.55,1.3);
  \draw[->](2.55,1.3)--(3,1.3);
  \draw[->](2,1.3)--(2,.1);
  \draw(.8,0)--(3,0);
  \draw[fill=white](1.15,1.2) rectangle(1.75,1.4);
  \draw[fill=white](-.1,.35) rectangle(.1,.95);
  \draw[fill=white](-.1+3,.35) rectangle(.1+3,.95);
  \draw[fill=white](-.1+2,.35) rectangle(.1+2,.95);
  \node[left] at(-.05,.65){$X$};
  \node[right] at(3.05,.65){$X$};
  \node[right] at(2.05,.65){$R$};
  \node at(.5,.6){\large$=$};
  \node[above] at (.95,1.3){$I_{1}$};
  \node[above] at (1.5,1.4){$r$};
  \node[above] at (2.5,1.3){$I_{2}$};
  \node[above] at (1.75,0){$I_{3}$};
\end{tikzpicture}}

\def\rysa{\begin{tikzpicture}
  \draw(0,0) rectangle(2,1.5);
  \draw(1,0) rectangle(3,2);
  \draw[white,fill=white](-.2,.65) rectangle(.2,.85);
  \draw(-.2,.65)--(.2,.65);
  \draw(-.2,.85)--(.2,.85);
  \draw[fill=white](.85,.6)--(1.15,.6)--(1,.9)--cycle;
  \draw(.85,.9)--(1.15,.9);
  \draw[fill=black](1,0) circle(.03) ;
  \draw[fill=black](1,1.5) circle(.03) ;
  \draw[white](2,.3) -- (2,1.2);
\draw [ snake=coil, segment amplitude=5pt, segment length=5pt]
(2,.3) -- (2,1.2);
\draw[white](3,.2)--(3,.6);
\draw[white](3,1.15)--(3,1.3);
\draw(2.9,1.15)--(3.1,1.15);
\draw(2.8,1.3)--(3.2,1.3);
  \draw[fill=black](3,.2) circle(.03) ;
  \draw[fill=black](3,.6) circle(.03) ;
  \draw(3,.6)--(3.2,.25);
  \node[below] at(3.3,.75){$K$};
  \node[below] at(1,0){$B$};
  \node[left] at(1.05,1.65){$A$};
\end{tikzpicture}}
\let\arcangle\angle
\long\def\nic#1{}


\setcounter{equation}0
\normalsize

\makeatletter
\def\@oddhead{\hbox to\hsize{\kern-7.5cm\rlap{\color{magenta!20}{\vrule height.9cm depth-1mm width22.5cm}}\kern2cm{\textcolor{white}{\LARGE\raise3mm\hbox{\textsf{\textbf{Liga zadaniowa Wydziału Matematyki, Informatyki i~Mechaniki,}}}}\hfill}\hfill\kern-3cm}}\makeatother

% LIGA--MATEMATYCZNA--ZDEFINIOWANIE
\ligaspis
 \long\def\matematyka
{ 
%\vspace{-30pt}
\advance\parskip by-2.5pt\baselineskip11.7pt plus.3pt minus.2pt 
\klub[2]{m}
  \Magenta{\textbf{Zadania z~matematyki nr 903, 904}}%\hfill{\scriptsize Termin nadsyłania rozwiązań: 30 VI 2025}%
  \marg[6]{Termin nadsyłania rozwiązań: \rlap{31 VIII 2025}}



\Black{\textit{Redaguje Marcin E. KUCZMA}}


{\bf903.}
Funkcje $f$~i~$g$ są dane wzorami
\begin{align*}
f(x)&=(1+x+x^2+x^3+\ldots+x^{42})+45x^{43},\\
g(x)&=(1+x+x^2+x^3+\;\;\ldots\;\;+x^{44})+45x^{45}.
\end{align*}
Liczby rzeczywiste $a,b$ spełniają warunki ${f(a)=g(b)=44}$. Wyjaśnić, która z~liczb $a,b$ jest większa.


{\bf904.}
Niech ${A_1=31}$, ${A_2=331}$, \dots, ${A_n=33\ldots31}$ ($n$~trójek, a~na końcu jedynka~-- zapis dziesiętny). Dla każdej liczby pierwszej ${p<44}$ rozstrzygnąć, czy istnieje liczba~$n$  (zależna od~$p$), dla której $A_n$ dzieli się przez~$p$. Jeśli istnieje, ustalić, czy jest wówczas nieskończenie wiele, czy tylko skończenie wiele takich numerów~$n$.

\textit{Zadanie 904 zaproponował  pan Witold Bednarek z~Łodzi.}
\marg[-40]{
    
    {\centering
Czołówka ligi zadaniowej {\bf Klub 44~M}\\
po uwzględnieniu ocen rozwiązań zadań\\[1pt]
 889 ($WT = 2{,}47$) i~890 ($WT = 1{,}77$)\\
z~numeru 11/2024\\
\vskip2pt}\tabcolsep5.5pt
\tabcolsep3pt
\begin{tabular}{@{}llr}
  Krzysztof Zygan&Lubin &42,03\\
Andrzej Kurach&Ryjewo &40,27\\
Andrzej Daniluk&Warszawa &37,89\\
Janusz Olszewski&Warszawa &37,17\\
Grzegorz Wiączkowski&   &36,58\\
Michał Warmuz&Żywiec &36,54\\
Marcin Kasperski&Warszawa &35,64\\
Krzysztof Kamiński&Pabianice &33,54\\
Marian Łupieżowiec&Gliwice &32,89\\
Jędrzej Biedrzycki&  &32,29\\
Marek Spychała&Warszawa &32,17\\
Krzysztof Maziarz&Londyn &31,96\\
  \end{tabular}\medskip

}

\Magenta{\textbf{Rozwiązania zadań z~numeru 2\slash2025}}

\Magenta{{\footnotesize Przypominamy treść zadań:}

\scriptsize
{\bf895.}
Na okręgu zaznaczono $n$ punktów; ${n\ge3}$ jest ustaloną liczbą nieparzystą. Każdemu punktowi została przyporządkowana wartość 0 lub~1. Dozwolone są ruchy polegające na wybraniu trzech kolejnych punktów o~wartościach (kolejno) $a,b,c$ takich, że ${a=c}$, i~zamianie $b$ na~$1{-}b$. Udowodnić, że startując z~dowolnej konfiguracji i~wykonując dozwolone ruchy można uzyskać jednakową wartość dla wszystkich $n$ punktów.


{\bf896.}
Udowodnić, że ciąg $(A_1,A_2,A_3,\ldots)$ o~wyrazach
$$
A_n=\frac1n{\textstyle\sum\limits_{k=1}^{2^n}}\frac1k
$$
jest malejący.


}\vspace*{\parskip}

\setlength{\columnsep}{.4cm}
\begin{dwieszpalty}\baselineskip11.7pt plus.3pt
{\bf895.}
Wygodnie będzie przejść na język kolorów: punkty z~wartością~0 --~białe; z~wartością~1 --~czarne. Układ $n$~punktów rozbijamy na bloki jednokolorowe. Co~najmniej jeden z~bloków ma długość nieparzystą; przyjmijmy (b.s.o.), że jest on czarny: 01111111110. Dozwolone ruchy polegają na zmianie koloru punktu (dowolnie wybranego), którego obaj sąsiedzi mają jednakowy kolor. Wykonując takie ruchy, przeprowadzamy podany przykładowo fragment do postaci 01010101010.

Kolejną serią ruchów pozbywamy się wszystkich punktów czarnych w~tym fragmencie: 00000000000. Skrajne zera, które na starcie obramowywały rozważany czarny blok, należały do jakichś bloków białych. Teraz one, wraz z~,,wybielonym'' fragmentem, zlewają się w~jeden dłuższy blok biały. Co istotne: liczba bloków zmniejszyła się. Powtarzamy procedurę (wybieramy blok jednobarwny długości nieparzystej, itd.); liczba bloków znów się zmniejszyła. Po dalszych powtórzeniach liczba bloków spadnie do~1. To będzie znaczyło, że wszystkie $n$~punktów uzyskało jednakowy kolor.



{\bf896.}
Niech ${B_n=H_n\slash\ln{n}}$, gdzie (jak zwykle) ${H_n=\sum_{k=1}^n\frac1k}$. Zauważmy, że ${B_{2^n}=A_n\slash{\ln2}}$. Wystarczy więc udowodnić, że ciąg $(B_n)$ jest malejący: ${B_{n-1}>B_n}$; czyli że
$$
H_n\ln(n-1)<H_{n-1}\ln{n}\quad\ \hbox{dla}\;\;n>1.
\leqno(1)
$$
Przydatne będą nierówności (dla ${n>1}$):
\begin{gather}
\frac1n<\ln{n}-\ln(n-1)<\frac1{n-1},
\tag2\\
\ln{n}<H_n.
\tag3
\end{gather}
Są one dobrze znane; krótkie uzasadnienie -- dla kompletności. Twierdzenie Lagrange'a dla funkcji $\;\ln{x}\;$ mówi, że \ ${\ln{n}-\ln(n{-}1)=1\slash{x_n}}$ \ dla pewnej liczby ${x_n\in(n{-}1,n)}$; stąd dwustronne oszacowanie~(2). Ustalmy ${m>1}$  i~wykonajmy sumowanie {\it prawej} części~(2) po ${n=2,\ldots,m}$, otrzymując zależność \ ${\ln{m}<H_{m{-}1}}$. \ Zastępujemy literkę~$m$ przez~$n$ i~mamy zależność~(3) (a~nawet nieco silniejszą, co tu bez znaczenia).

Teraz krótki rachunek; w~pierwszym kroku wykorzystujemy {\it lewą} część~(2), a~w~ostatnim -- zależność~(3):\vspace*{-3pt}
\begin{align*}
H_n\ln(n-1)&<H_n\biggl(\ln{n}-\frac1n\biggr)\\&=\biggl(H_{n-1}+\frac1n\biggr)\ln{n}+H_n\biggl(-\frac1n\biggr)={}\cr
&=H_{n-1}\ln{n}+\frac1n\biggl(\ln{n}-H_n\biggr)<H_{n-1}\ln{n}
\end{align*}
-- i~mamy tezę~(1).
\end{dwieszpalty}


}



% -------------------- LIGA--FIZYCZNA--ZDEFINIOWANIE ---------------------------------------

\long\def\fizyka
{\klub[5]{f}
\Magenta{\textbf{Zadania z~fizyki nr 800, 801}}%
\marg{Termin nadsyłania rozwiązań: \rlap{31 VIII 2025}\medskip


}

\Black{\textit{Redaguje Elżbieta ZAWISTOWSKA}}%

\textbf{800.} W~obwodzie przedstawionym na rysunku 1 dioda jest idealna i~wszystkie opory omowe zaniedbywalne. Klucz $K$ zamknięto na czas $\tau$, a~następnie otwarto. W~chwili otwierania klucza natężenie prądu w~cewce miało wartość $I_{0}$. Po jakim czasie po otwarciu klucza osiągnęło  ono wartość maksymalną równą  $2I_{0}$? Narysuj wykres zależności natężenia prądu w~cewce od czasu $t$ ($0<t<\infty$).\marg{Rys. 1\rysa%\\[4pt]Rys. 1

}

\textbf{801.} Statek kosmiczny przemieszcza się od Ziemi do Marsa po orbicie eliptycznej, której peryhelium znajduje się na orbicie Ziemi, a~aphelium na orbicie Marsa. Znaleźć czas przelotu oraz minimalny czas, w~jakim kosmonauci będą musieli oczekiwać na Marsie na start w~drogę powrotną po takiej samej orbicie. Okres obiegu Ziemi wokół Słońca wynosi $T_{\rm Z}=365{,}25$~dób, okres obiegu Marsa  $T_{\rm M}=687$~dób. Przyjąć, że planety poruszają się po orbitach kołowych leżących w~jednej płaszczyźnie.
\marg{
  
\bigskip
{{{\centering
Czołówka ligi zadaniowej {\bf Klub 44~F}\\
po uwzględnieniu ocen rozwiązań zadań\\
786 ($WT=2{,}5$), 787 ($WT=3{,}25$)\\  z~numeru 11/2024
\vskip2pt}
\spaceskip.25em\tabcolsep3pt\def\pp(##1){&##1}
\begin{tabular}{@{}l@{\ }@{\!}r}
%  Konrad Kapcia (Poznań)&                         \Magenta{3}\,–\,\Magenta{44}+3,25\\
%Paweł Perkowski (Ożarów Maz.)&          \Magenta{6}\,–\,\Magenta{44}+0,70\\
%Jacek Konieczny (Poznań) &                           40,87\\
Tomasz Wietecha (Tarnów)&                17\,–\,42,90\\
Jan Zambrzycki (Białystok)&                    4\,–\,29,64\\
Andrzej Nowogrodzki\\\kern2.4cm (Chocianów)&     3\,–\,27,49\\
\end{tabular}\bigskip}}



}

\Magenta{\textbf{Rozwiązania zadań z~numeru 2\slash2025}

\footnotesize{Przypominamy treść zadań:}
 
{\scriptsize 

{\bf 792.}  $N=100$ jednakowych kulek o~średnicy $d=0{,}1$~mm znajduje się w~ustawionym pionowo cylindrycznym naczyniu o~podstawie $S=1$~m$^{2}$ pod nieruchomym tłokiem, który znajduje się na wysokości $h=1$~m. Kulki poruszają się chaotycznie ze średnią prędkością kwadratową $v_{0}=100$~m/s. Tłok zaczęto podnosić z~prędkością $u=1$~m/s i~został on zatrzymany na wysokości~$2h$. Jaka średnia prędkość kulek ustaliła się po długim czasie? Nie ma strat energii mechanicznej podczas zderzeń, nie uwzględniamy siły grawitacji.
                                                                                                                                                                                                                                                                                                                                                                                                                                                                                                                                                                                 
{\bf 793.} Do źródła o~sile elektromotorycznej $U=1{,}5$~V i~zaniedbywalnym oporze wewnętrznym dołączono długi łańcuch jednakowych amperomierzy i~taką samą liczbę jednakowych woltomierzy (rys.~2). Opór wewnętrzny amperomierza wynosi $r=1\ \Omega$, woltomierza 10~k$\Omega$. Jakie są wskazania pierwszego i~drugiego amperomierza? Ile wynosi suma wskazań wszystkich amperomierzy oraz suma wskazań wszystkich woltomierzy w~łańcuchu?

}}
 \vspace*{-\parskip}

{\scriptsize Rys. 2 \lfb%\\[4pt]Rys. 2
\kern1cm  \qquad Rys. 3  \rysc%\includegraphics[width=4cm]{lf-r793}%\\[4pt]Rys. 3

}

 \vspace*{\parskip}
\begin{dwieszpalty}\baselineskip11.7pt plus.3pt
\textbf{792.}
Średnia droga swobodna kulki w~naczyniu, czyli droga przebyta przez nią
do chwili zderzenia z~inną kulką, wynosi
\(\lambda = 1/( \pi d^{2}n )\), gdzie \(n = N/(Sh)\) jest
liczbą kulek w~jednostce objętości. Czas między dwoma kolejnymi
zderzeniami \(\tau = \xfrac{\lambda}{v_{0}} = 3 \cdot 10^{3}{\rm s} \gg \xfrac{h}{v_{0}} = 10^{- 2}\)s,
zatem kulki praktycznie nie zderzają się między sobą. Podczas ruchu
tłoka zmiana energii kulki związana jest tylko z~uderzeniami o~tłok i~określona przez zmianę jej pionowej składowej prędkości.

Niech odległość tłoka od dna naczynia w~pewnej chwili wynosi \(l\), tłok
porusza się z~prędkością \(u\), a~kulka dogania go z~prędkością
\(v \gg u\). W~układzie związanym z~tłokiem prędkość kulki wynosi
\(v - u\), natomiast po odbiciu \(- v + u\). W~układzie cylindra
prędkość kulki po odbiciu to \(- v + 2u\), czyli jej wartość zmniejsza
się o \(2u\). Kolejne uderzenie o~tłok nastąpi po czasie
\(\mathrm{\Delta}t \cong \xfrac{2l}{v}\). Tłok przesunie się w~tym czasie
o
\(\mathrm{\Delta}l = u\mathrm{\Delta}t = - l\xfrac{\mathrm{\Delta}v}{v}\),
stąd\vspace*{-4pt}
\[v\mathrm{\Delta}l + l\mathrm{\Delta}v = 0,\]
co oznacza, że iloczyn \(lv\) pozostaje stały.

Dla składowej pionowej prędkości kulki mamy
\[hv_{y0} = 2hv_{y},\qquad v_{y} = \xfrac{v_{y0}}{2,}\]
czyli na wysokości \(2h\) pionowa składowa prędkości kulki jest dwa razy
mniejsza niż na wysokości \(h\).

Początkowa energia kinetyczna kulki wynosiła
\[W_{0} = \xfrac{M( {v_{x0}}^{2} + {v_{y0}}^{2} + {v_{z0}}^{2} )}{2},\]
 gdzie 
\(( {v_{x0}}^{2} )_{\rm śr} = ( {v_{y0}}^{2} )_{\rm śr} = ( {v_{z0}}^{2} )_{\rm śr}\).

Po przemieszczeniu tłoka zmieniła się tylko składowa pionowa prędkości
kulki, i~jej prędkość końcowa wynosi:
\[W = M\xfrac{( {v_{x0}}^{2} + \xfrac{{v_{y0}}^{2}}{4} + {v_{z0}}^{2} )}{2}.\]
Po długim czasie, gdy ruch kulek ponownie staje się w~pełni chaotyczny,
średnia wartość ich prędkości wynosi:
\[v_{\rm śr} = v_{0}\sqrt{\xfrac{3}{4}} = 87\ \rm\xfrac{m}{s.}\]
\textbf{793.}
a) Prąd płynący przez pierwszy amperomierz jest równy
\(I_{1} = \xfrac{U}{X}\), gdzie \(X\) to opór zastępczy układu.
Po dołączeniu z~lewej strony jeszcze jednego oczka (rys.~3) opór
zastępczy nie zmienia się.
Stąd
\(X = r + \xfrac{RX}{(R + X)}\), \(X = \xfrac{( r + \sqrt{r^{2} + 4Rr} )}{2 = 100{,}5\ \Omega}\),
\(I_{1} = 14{,}9\ \rm mA\).

b) Część obwodu bez pierwszego amperomierza ma opór
\(X - r = 99{,}5\ \Omega\).
Natężenie prądu płynącego przez drugi amperomierz
\(I_{2} = \xfrac{U_{2}}{X}\), gdzie
\(U_{2} = U - rI_{1} = U( 1 - \xfrac{r}{X} ) = 0{,}99U\),
\(I_{2} = 0{,}99I_{1} = 14,8\ \rm mA\).

c) Suma napięć na wszystkich amperomierzach równa jest napięciu ogniwa
\(U = 1,5\ V\). Prąd płynący przez amperomierz to stosunek jego napięcia
i~oporu, suma natężeń prądów na wszystkich amperomierzach wynosi więc
\(I = \xfrac{U}{r} = 1,5\ \rm A\).

d) Suma natężeń prądów płynących przez wszystkie woltomierze równa jest
natężeniu prądu
płynącego przez pierwszy amperomierz, stąd suma napięć na wszystkich
woltomierzach
wynosi \(RI_{1} = 149\ V\).


\end{dwieszpalty}


}





\makeatletter
%\def\@oddhead{\ifnum\thepage=17\vbox{\hbox to\hsize{\kern-7.5cm\rlap{\color{magenta!20}{\vrule height.9cm depth-1mm width22.5cm}}\kern2cm{\textcolor{white}{\LARGE\raise3mm\hbox{\textsf{\textbf{Liga zadaniowa Wydziału Matematyki, Informatyki i~Mechaniki,}}}}\hfill}\hfill\kern-3cm}\vspace{-4pt}
%  \hbox to\hsize{\kern-7.5cm\rlap{\color{magenta!20}{\vrule height.9cm depth-1mm width22.5cm}}\kern2cm{\textcolor{white}{\LARGE\raise3mm\hbox{\textsf{\textbf{Wydziału Fizyki Uniwersytetu Warszawskiego i~Redakcji \textit{Delty}}}}}\hfill}\hfill\kern-3cm}}\else\ifodd\thepage\hbox to\hsize{\kern-7.5cm\rlap{\color{magenta!20}{\vrule height.9cm depth-1mm width22.5cm}}\kern2cm{\textcolor{white}{\LARGE\raise3mm\hbox{\textsf{\textbf{Wydziału Fizyki Uniwersytetu Warszawskiego i~Redakcji \textit{Delty}}}}}\hfill}\hfill\kern-3cm}\else\hbox to\hsize{\kern-7.5cm\rlap{\color{magenta!20}{\vrule height.9cm depth-1mm width22.5cm}}\kern2cm{\textcolor{white}{\LARGE\raise3mm\hbox{\textsf{\textbf{Liga zadaniowa Wydziału Matematyki, Informatyki i~Mechaniki,}}}}\hfill}\hfill\kern-3cm}\fi\fi}\makeatother
  
\normalsize

\vspace*{-20pt}
{\fizyka}
\eject
\makeatletter
\def\@oddhead{\hbox to\hsize{\kern-7.5cm\rlap{\color{magenta!20}{\vrule height.9cm depth-1mm width22.5cm}}\kern2cm{\textcolor{white}{\LARGE\raise3mm\hbox{\textsf{\textbf{Wydziału Fizyki Uniwersytetu Warszawskiego i~Redakcji \textit{Delty}}}}}\hfill}\hfill\kern-3cm}}\makeatother



{\matematyka}
\vfill\vskip-4pt
\regulamin



\normalsize
\eject\setcounter{equation}0
\makeatletter
\def\@oddhead{\hfill}
\makeatother
\normalsize

